Four lines of work bear on the question of which units a corpus should be indexed
by, and where the decision comes from.

\paragraph{Multiple representations per item.}
Giving a document several embeddings rather than one is well established.
Multi-view document representation learns several latent viewer embeddings per
document with a global-local objective~\cite{zhang2022multiview}; multi-aspect
dense retrieval learns explicit aspect embeddings from product
metadata~\cite{kong2022madral}; MC-indexing indexes each section under three
fixed views, raw text, keywords and summary~\cite{dong2024mcindexing}; and late
interaction is the limiting case where views are tokens~\cite{khattab2020colbert}.
Vector databases now ship the mechanism directly as named vectors with per-view
indexes. In every case the view set is fixed in advance, by a designer, by a
learned architecture, or by available metadata. None varies with the question the
index is being built to answer, and our results show that this is the property
that decides whether views help at all.

\paragraph{Choosing the unit of indexing.}
A parallel line asks what a retrieval unit should be. Proposition indexing
rewrites documents into atomic propositions~\cite{chen2024densex}; contextual
retrieval prepends situating text to each chunk before
embedding~\cite{anthropic2024contextual}; LumberChunker and its successors place
chunk boundaries with a model~\cite{duarte2024lumberchunker}; and
mixture-of-granularity routes queries among pre-segmented
granularities~\cite{zhong2024mog}. These decide a granularity, globally and once,
and several require the model to read every item, which is the cost this work
avoids. They do not produce semantically named views, and the decision is not
conditioned on a stated goal.

\paragraph{Generative models that build index structures.}
GraphRAG has a model extract an entity graph and community summaries for
query-focused sensemaking~\cite{edge2024graphrag}, with auto prompt tuning
generating the extraction schema from roughly a one percent corpus
sample~\cite{microsoft2024autotune}; RAPTOR builds a recursive summary
tree~\cite{sarthi2024raptor}; EVAPORATE has a model synthesise extraction
functions from a sample which then run over the whole corpus, reducing token cost
by two orders of magnitude~\cite{arora2023evaporate}; and StructRAG selects a
corpus structure type per task from a fixed menu~\cite{li2024structrag}. The
closest of all is TnT-LLM, which derives a label taxonomy from corpus samples
under a use-case instruction and distils it into lightweight classifiers that
execute at scale~\cite{wan2024tntllm}: the same economics as our design phase,
producing a flat taxonomy for classification rather than a set of embedding views
serving retrieval. Most recently CAMI has an agent propose candidate enrichment
indexes and select a portfolio under a cost budget~\cite{qidwai2026cami}, which
conditions index construction on the corpus and a budget but not on what the user
is trying to find out.

\paragraph{Conditioning representations on instructions.}
Instruction-following encoders produce a different vector for the same text under
a different instruction~\cite{su2023instructor,weller2024followir,weller2025promptriever},
and recent work materialises per-instruction views cheaply without re-encoding
the corpus~\cite{sun2025gstransform}. These subsume the execution layer of our
method: given $k$ instructions, they produce $k$ views. What they do not supply is
the step above, deciding which $k$ views a corpus and an objective call for. We
therefore treat an instruction-conditioned encoder given the objective text as the
primary baseline rather than as related work, since it is the cheapest credible
alternative to the entire approach.

\paragraph{Position of this work.}
Against that background we make no claim to multiple views per item, to generative
design from a sample, or to agentic index selection. We ask a narrower question
that, to our knowledge, none of the above tests: does conditioning the view set on
a free-form statement of the research objective change which views are chosen, and
does that change the retrieval quality of the resulting index? We answer it with a
conditioning ladder in which every rung shares an encoder, a builder and an index,
and only what the designer is shown varies, and with a cross-over design in which
two objectives over one corpus must each be better served by their own view set.
